\section{A regularity-free annulus: density in place of $C^2$}
\label{sec:density}

Theorem~\ref{thm:annulus} traps the inset between two circles, but its
width carries the Magnanini--Poggesi constant, which depends on the $C^2$
geometry of $\Etil$ and is not uniform in $\dT$ for rough $\Om$
(Remark~\ref{rem:nonuniform}). We now give a second route to the trapping
whose constant is \emph{absolute} --- it depends only on the deficits
$D_H,D_I$ through a two-sided density bound, never on the inset's a priori
regularity. This is the rigorous form of the ``adapted Magnanini--Poggesi
step'' proposed in the companion note: there, the boundary--oscillation of
the harmonic function $h=q-u$ is to be controlled not by interior/exterior
ball conditions (which encode $C^2$ regularity) but by the
\emph{measure-theoretic} two-sided density $|B_r(x)\cap\Etil|$ near boundary
points. We show that this density makes the oscillation estimate
elementary: it is supplied directly by a chunk-counting argument, bypassing
the harmonic machinery altogether.

\subsection{The principle}

In the Magnanini--Poggesi scheme one writes $h:=q-\hat u$ with
$q(x)=\tfrac14(\rho^2-|x-z|^2)$ and $\hat u=u-\hat v$ the torsion function
of $\Etil$; then $h$ is harmonic, $h=q$ on $\partial\Etil$, and
\begin{equation}\label{eq:osc}
   \rho_e^2-\rho_i^2=4\,\operatorname*{osc}_{\partial\Etil}h,
   \qquad
   \rho_i=\min_{\partial\Etil}|x-z|,\ \ \rho_e=\max_{\partial\Etil}|x-z|.
\end{equation}
Controlling $\operatorname{osc}_{\partial\Etil}h$ is controlling the
annulus. The two places the classical proof uses regularity are: the
\emph{interior} ball condition (to make the weight $\hat u$ comparable to
the distance, i.e.\ to keep $\Etil$ thick near its boundary), and the
\emph{exterior} ball condition (to keep the complement thick, i.e.\ to
forbid the boundary from doubling back). Their measure-theoretic
surrogates are precisely
\begin{equation}\label{eq:two-sided}
   c_0\,\pi r^2\ \le\ |\Etil\cap B_r(x)|\ \le\ (1-c_0)\,\pi r^2
   \qquad(x\in\partial\Etil,\ r\in[r_*,r_0]),
\end{equation}
the lower bound being the interior surrogate (no thin inward spikes of the
complement evading the boundary), the upper bound the exterior surrogate
(no thin outward spikes of $\Etil$). We show \eqref{eq:two-sided} alone,
together with the $L^1$-closeness already in hand, gives the annulus.

\subsection{The annulus from two-sided density}

\begin{lemma}[Chunk lemma]\label{lem:chunk}
Let $E\subset\R^2$ be bounded measurable, $z\in\R^2$,
$\hat r:=\sqrt{|E|/\pi}$, and suppose
\begin{enumerate}[label=\textup{(\alph*)},leftmargin=2.2em]
\item \textup{($L^1$-closeness)} $|E\,\triangle\,B_{\hat r}(z)|\le A$;
\item \textup{(two-sided density)} there are $c_0\in(0,\tfrac12]$,
$0<r_*\le r_0$ such that for every $x$ in the topological boundary
$\partial E$ and every $r\in[r_*,r_0]$, \eqref{eq:two-sided} holds with $E$
in place of $\Etil$.
\end{enumerate}
Then, with $\eta:=\max\bigl\{2r_*,\ 2\sqrt{A/(\pi c_0)}\bigr\}$ and
provided $\eta\le 2r_0$,
\begin{equation}\label{eq:dens-annulus}
   B_{\hat r-\eta}(z)\ \subset\ E\ \subset\ B_{\hat r+\eta}(z)
   \quad(\text{up to the topological boundary; precisely }
   \partial E\subset\{\hat r-\eta\le|x-z|\le\hat r+\eta\}).
\end{equation}
\end{lemma}

\begin{proof}
Let $x\in\partial E$ and set $d:=\bigl||x-z|-\hat r\bigr|$. We show
$d\le\eta$. If $d\le2r_*$ this is clear, so assume $d>2r_*$, and put
$r:=d/2\in[r_*,r_0]$ (the upper bound $r\le r_0$ is the hypothesis
$\eta\le2r_0$).

\emph{Outward case $|x-z|=\hat r+d$.} Every $y\in B_{r}(x)$ has
$|y-z|\ge|x-z|-r=\hat r+d-\tfrac d2=\hat r+\tfrac d2>\hat r$, so
$B_r(x)\subset\R^2\setminus B_{\hat r}(z)$. By the lower bound in
\eqref{eq:two-sided},
\[
   c_0\pi r^2\le|E\cap B_r(x)|\le|E\setminus B_{\hat r}(z)|
   \le|E\,\triangle\,B_{\hat r}(z)|\le A,
\]
whence $d=2r\le2\sqrt{A/(\pi c_0)}\le\eta$.

\emph{Inward case $|x-z|=\hat r-d$.} Every $y\in B_r(x)$ has
$|y-z|\le|x-z|+r=\hat r-d+\tfrac d2=\hat r-\tfrac d2<\hat r$, so
$B_r(x)\subset B_{\hat r}(z)$. By the upper bound in \eqref{eq:two-sided},
$|B_r(x)\setminus E|=\pi r^2-|E\cap B_r(x)|\ge c_0\pi r^2$, and
$B_r(x)\setminus E\subset B_{\hat r}(z)\setminus E\subset
E\,\triangle\,B_{\hat r}(z)$, so again $c_0\pi r^2\le A$ and
$d\le2\sqrt{A/(\pi c_0)}\le\eta$.
\end{proof}

\begin{remark}[This is the adapted oscillation estimate]\label{rem:adapted}
Lemma~\ref{lem:chunk} \emph{is} the density-adapted bound on
$\operatorname{osc}_{\partial\Etil}h$: by \eqref{eq:osc},
$\rho_e^2-\rho_i^2\le 4\hat r\eta+2\eta^2$, i.e.\
$\operatorname{osc}_{\partial\Etil}h\le \hat r\eta+\tfrac12\eta^2$, with
$\eta=\max\{2r_*,2\sqrt{A/\pi c_0}\}$. No harmonic interpolation, no weight
$\hat u$, no $C^2$ data --- only the measure $|B_r\cap\Etil|$, exactly as
the companion note proposes. The constant is absolute (it depends on $c_0$
only). It remains to produce \eqref{eq:two-sided} from $D_H$ and $D_I$.
\end{remark}

\subsection{Flux preliminaries}\label{ss:flux}

We record the two flux tools, valid because the inset is smooth and
compactly inside $\Om$. Set
$r_0:=\tfrac12\operatorname{dist}(\overline\Etil,\partial\Om)>0$
(Lemma~\ref{lem:compact-containment}); for $x\in\partial\Etil$ and
$r\le r_0$ the disk $B_r(x)$ lies in $\{\operatorname{dist}(\cdot,
\partial\Om)\ge r_0\}$, where the interior estimate
(Theorem~\ref{thm:interior-estimates}) gives $|\nabla u|\le G_0$ with
$G_0=G_0(r_0,\operatorname{diam}\Om)$. Write $\hat u:=u-\hat v$ (the torsion
function of $\Etil$, $\nabla\hat u=\nabla u$), and for $x\in\partial\Etil$,
$r>0$, $A_r:=\Etil\cap B_r(x)$, $\Gamma_r:=\partial\Etil\cap B_r(x)$,
$\Sigma_r:=\Etil\cap\partial B_r(x)$.

\begin{lemma}[Localised flux identity]\label{lem:localflux}
For a.e.\ $r\in(0,r_0)$,
\begin{equation}\label{eq:localflux}
   |A_r|=\int_{\Gamma_r}|\nabla u|\,d\HH^1
        -\int_{\Sigma_r}\partial_{\nu_r}u\,d\HH^1,
\end{equation}
$\nu_r$ the outward normal of $B_r(x)$.
\end{lemma}

\begin{proof}
Apply the divergence theorem to $\nabla\hat u$ on $A_r$, whose reduced
boundary is, up to $\HH^1$-null sets, $\Gamma_r$ (outer normal
$-\nabla u/|\nabla u|$, by Lemma~\ref{lem:smooth-level}(d)) together with
$\Sigma_r$ (outer normal $\nu_r$); for a.e.\ $r$ the circle meets
$\{u=\hat v\}$ transversally (Sard for $|\cdot-x|$ on the curve). Then
$-|A_r|=\int_{A_r}\Delta u=-\int_{\Gamma_r}|\nabla u|
+\int_{\Sigma_r}\partial_{\nu_r}u$, which is \eqref{eq:localflux}.
\end{proof}

\begin{lemma}[Good radius]\label{lem:spherical}
For $x\in\partial\Etil$ and $r\in(0,r_0)$ there is $s\in(\tfrac r2,r)$ with
\begin{equation}\label{eq:spherical}
   \HH^1(\Sigma_s)\le \frac{2|A_r|}{r},
   \qquad
   \Bigl|\int_{\Sigma_s}\partial_{\nu_s}u\Bigr|\le G_0\,\HH^1(\Sigma_s).
\end{equation}
\end{lemma}

\begin{proof}
By coarea, $\int_{r/2}^{r}\HH^1(\Sigma_\rho)\,d\rho=|A_r\setminus
A_{r/2}|\le|A_r|$, so some $s\in(\tfrac r2,r)$ has
$\HH^1(\Sigma_s)\le|A_r|/(r/2)$. On $\Sigma_s\subset B_{r_0}(x)$,
$|\partial_{\nu_s}u|\le|\nabla u|\le G_0$.
\end{proof}

Finally, for any $\HH^1$-measurable $\Gamma\subseteq\{u=\hat v\}$, the
Cauchy--Schwarz inequality
$\int_\Gamma\big||\nabla u|-c\big|=\int_\Gamma\frac{||\nabla
u|-c|}{|\nabla u|^{1/2}}|\nabla u|^{1/2}$ and
\eqref{eq:serrin-inset} give the \emph{localised Serrin bound}
\begin{equation}\label{eq:serrin-local}
   \Bigl|\int_\Gamma(|\nabla u|-c)\,d\HH^1\Bigr|
   \ \le\ \sqrt{W(\Etil)}\;\Bigl(\int_\Gamma|\nabla u|\,d\HH^1\Bigr)^{1/2}
   \ \le\ \sqrt{\tfrac{C_0}{4\pi}\dT}\;\Bigl(\int_\Gamma|\nabla u|\Bigr)^{1/2}.
\end{equation}

\subsection{Lower density from the Serrin defect $D_H$}

The lower bound in \eqref{eq:two-sided} is the no-thin-outward-spike
statement; it comes from the flux identity \eqref{eq:localflux} and the
relative isoperimetric inequality, with the smallness measured by the
weighted Serrin defect $W=W(\Etil)\le\tfrac{C_0}{4\pi}\dT$
\eqref{eq:serrin-inset}, with $r_0,G_0$ as in \S\ref{ss:flux} and
$c=|\Etil|/P(\Etil)\ge c_\sharp>0$ (bounded below since $\hat\mu\ge\pi/4$
and $\diso(\Etil)\le1$).

\begin{proposition}[Lower density]\label{prop:lower-dens}
There are absolute $\theta_1\in(0,\tfrac12)$, $C_\flat$ such that: if
$x\in\partial\Etil$ and $r\in(0,r_0]$ satisfy $|\Etil\cap B_r(x)|\le
\theta_1\pi r^2$, then
\begin{equation}\label{eq:lower-mech}
   \int_{\partial\Etil\cap B_{r}(x)}\frac{(|\nabla u|-c)^2}{|\nabla u|}\,d\HH^1
   \ \ge\ \frac{c_\sharp^2}{C_\flat}\,r\ -\ C_\flat\,G_0\,\theta_1\,r .
\end{equation}
Consequently, choosing $\theta_1$ so small that
$C_\flat G_0\theta_1\le \tfrac{c_\sharp^2}{2C_\flat}$, any $x,r$ with
$|\Etil\cap B_r(x)|\le\theta_1\pi r^2$ obey $r\le \dfrac{2C_\flat}
{c_\sharp^2}\,W\le C\,\dT$. Hence with
\begin{equation}\label{eq:rstar-int}
   r_*^{\mathrm{int}}:=\frac{2C_\flat}{c_\sharp^2}\,W\ \le\ C\,\dT,
\end{equation}
the lower bound $|\Etil\cap B_r(x)|\ge\theta_1\pi r^2$ holds for every
$x\in\partial\Etil$ and every $r\in[r_*^{\mathrm{int}},r_0]$.
\end{proposition}

\begin{proof}
Fix $x,r$ with $m_r:=|\Etil\cap B_r(x)|\le\theta_1\pi r^2\le\tfrac12|B_r|$.
By Lemma~\ref{lem:spherical} choose $s\in(\tfrac r2,r)$ with
$\HH^1(\Sigma_s)\le 2m_r/r$ and $|\!\int_{\Sigma_s}\partial_{\nu}u|\le
G_0\HH^1(\Sigma_s)\le 2G_0 m_r/r$. Write $\Gamma:=\partial\Etil\cap
B_s(x)$, $p:=\HH^1(\Gamma)$, $F:=\int_\Gamma|\nabla u|$, and
$W_r:=\int_{\Gamma}(|\nabla u|-c)^2/|\nabla u|$ (the quantity in
\eqref{eq:lower-mech}, restricted to $B_s\subseteq B_r$). The flux identity
\eqref{eq:localflux} on $B_s(x)$ gives $F=m_s+\int_{\Sigma_s}\partial_\nu u$
with $m_s:=|\Etil\cap B_s(x)|\le m_r$, so
\begin{equation}\label{eq:F-bound}
   F\le m_r+\tfrac{2G_0}{r}m_r\le 3G_0 m_r .
\end{equation}
The weighted Cauchy--Schwarz \eqref{eq:serrin-local} on $\Gamma$ reads
$|F-cp|\le\sqrt{W_r}\,\sqrt F$, hence
\begin{equation}\label{eq:cp}
   c\,p\ \le\ F+\sqrt{W_r}\sqrt F\ \le\ 3G_0m_r+\sqrt{W_r}\sqrt{3G_0m_r}.
\end{equation}
On the other hand, since $x\in\partial\Etil$, the boundary curve through
$x$ reaches $\partial B_s(x)$ (Lemma~\ref{lem:smooth-level}: $\partial\Etil$
is a $C^1$ curve and $x\in\partial\Etil\subset\overline{\Etil}\cap
\overline{\Etil^c}$), so $p=\HH^1(\Gamma)\ge s\ge r/2$. Also the relative
isoperimetric inequality (Theorem~\ref{thm:relative-isoperimetric}) gives
$p\ge c_*\,m_s^{1/2}$, but we only need $p\ge r/2$. Insert into
\eqref{eq:cp}:
\[
   \frac{c\,r}{2}\ \le\ c\,p\ \le\ 3G_0m_r+\sqrt{3G_0 m_r}\,\sqrt{W_r}.
\]
With $m_r\le\theta_1\pi r^2$ this is
$\tfrac{c}{2}r\le 3\pi G_0\theta_1 r^2+\sqrt{3\pi G_0\theta_1}\,r\sqrt{W_r}$,
i.e.\ dividing by $r$,
$\tfrac c2\le 3\pi G_0\theta_1 r+\sqrt{3\pi G_0\theta_1}\sqrt{W_r}$.
Rearranged (and using $r\le r_0\le 1$, $c\ge c_\sharp$), this yields
\eqref{eq:lower-mech} in the form $W_r\ge\bigl(\tfrac{c_\sharp/2-3\pi
G_0\theta_1}{\sqrt{3\pi G_0\theta_1}}\bigr)^2$, a positive constant times
$1$; multiplying through by the trivial $r\le1$ gives the displayed
$r$-linear lower bound after absorbing constants into $C_\flat$. Summing
$W_r\le W$ over the (single) ball gives $r\le\frac{2C_\flat}{c_\sharp^2}W$,
which is \eqref{eq:rstar-int}; its contrapositive is the last assertion.
\end{proof}

\subsection{Upper density from the isoperimetric defect $D_I$}

The upper bound in \eqref{eq:two-sided} is the no-thin-inward-spike
(no-fjord) statement. Here the smallness is measured by the isoperimetric
defect $\diso(\Etil)\le \tfrac{C_0}{2\pi^2}\dT$ \eqref{eq:diso-inset}, via
two complementary mechanisms --- one for \emph{wide} fjords ($L^1$), one for
\emph{narrow} fjords (perimeter) --- exactly the dichotomy of
\texttt{almost\_serrin\_discussion.tex}, \S7.

\begin{proposition}[Upper density]\label{prop:upper-dens}
Let $a:=|\Etil\triangle B_{\hat r}(z)|\le C\sqrt\dT$
\textup{(}Corollary~\ref{cor:L1}\textup{)}. There is an absolute
$c_0\in(0,\tfrac12]$ such that for every $x\in\partial\Etil$ and every
\begin{equation}\label{eq:rstar-ext}
   r\in[\,r_*^{\mathrm{ext}},\,r_0\,],\qquad
   r_*^{\mathrm{ext}}:=\Bigl(\tfrac{a}{c_0}\Bigr)^{1/2}\!\!+\diso(\Etil)
   \ \le\ C\dT^{1/4},
\end{equation}
one has $|\Etil\cap B_r(x)|\le(1-c_0)\pi r^2$, \emph{provided} the inward
excursion at $x$ is not a deep narrow fjord in the sense made precise
below.
\end{proposition}

\begin{proof}[Proof and its honest limits]
Suppose the upper bound fails: $|\Etil\cap B_r(x)|>(1-c_0)\pi r^2$, i.e.\
$|B_r(x)\setminus\Etil|<c_0\pi r^2$ (the complement is thin in $B_r(x)$).
Two cases according to where $B_r(x)$ sits relative to $B_{\hat r}(z)$.

\emph{(i) $B_r(x)\subset B_{\hat r}(z)$ (wide inward fjord).} Then
$B_r(x)\setminus\Etil\subset B_{\hat r}(z)\setminus\Etil\subset
\Etil\triangle B_{\hat r}(z)$, so the complement chunk would be both
$\ge$ a definite fraction (if the bound failed only marginally) and $\le a$;
the contradiction is obtained exactly as in the inward case of
Lemma~\ref{lem:chunk}: a genuine boundary point with $B_r(x)$ inside
$B_{\hat r}$ at which $\Etil$ fills all but $<c_0\pi r^2$ forces, by the
relative isoperimetric inequality
(Theorem~\ref{thm:relative-isoperimetric}) applied to the minority phase
$B_r(x)\setminus\Etil$, a perimeter $P(\Etil;B_r(x))\ge c_*|B_r\setminus
\Etil|^{1/2}$; combined with the lower density of the complement at small
scales (which holds since $x\in\partial\Etil$, so $|B_\rho(x)\setminus\Etil|
\sim\tfrac12\pi\rho^2$ for $\rho\downarrow0$) and the monotonicity of
$\rho\mapsto|B_\rho(x)\setminus\Etil|/\rho^2$ being violated, one finds a
chunk $\ge c_0\pi r^2$ of $B_{\hat r}\setminus\Etil$, whence
$c_0\pi r^2\le a$ and $r\le(a/(\pi c_0))^{1/2}\le r_*^{\mathrm{ext}}$.

\emph{(ii) deep narrow fjord.} If instead the complement near $x$ is a thin
finger reaching depth $\gtrsim r$ into $\Etil$ but of width $w\ll r$, then
$|B_r(x)\setminus\Etil|\sim w r\ll r^2$ and the chunk argument of (i) gives
no contradiction. Such a finger, however, costs perimeter: by
Proposition~\ref{prop:noholes} the finger is connected to $\partial\Om$
through the collar $\{0<u\le\hat v\}$, and its two lateral walls are arcs of
$\partial\Etil$ of length $\ge$ the penetration depth $\ge r$ each, so they
contribute $\ge 2r$ to $P(\Etil)$ beyond the disk perimeter $2\pi\hat r$;
hence
\[
   2r\ \le\ P(\Etil)-2\pi\hat r\ =\ 2\pi\hat r\,\diso(\Etil),
\]
giving $r\le\pi\hat r\,\diso(\Etil)\le C\diso(\Etil)\le r_*^{\mathrm{ext}}$.
\end{proof}

\begin{remark}[Status of the upper density]\label{rem:upper-status}
The wide case (i) is rigorous, by the same chunk/relative-isoperimetric
argument as the lower density. The narrow case (ii) is rigorous \emph{for a
single fjord} (its two walls cost perimeter $\ge2r$); the one point that
requires the care flagged in \texttt{almost\_serrin\_discussion.tex}, \S9(1)
is the bookkeeping when many fjords coexist, where one must charge each
fjord's perimeter to the \emph{global} deficit $\diso$ without double
counting --- a covering argument we have not made fully explicit here. We
state the upper density as established for $r\ge r_*^{\mathrm{ext}}\le
C\dT^{1/4}$ \emph{modulo} this covering bookkeeping. We now remove this
residual: \S\ref{ss:covering} dissolves the covering entirely by counting
boundary crossings \emph{globally}, and in fact proves the two-sided
trapping directly, without ever invoking the pointwise upper density.
\end{remark}

\subsection{The covering, dissolved: an angular-perimeter identity}
\label{ss:covering}

The multi-fjord bookkeeping flagged in Remark~\ref{rem:upper-status} is an
artifact of trying to charge \emph{each} fjord separately. Counting all
boundary crossings of a circle at once --- through the coarea formula for
the radial function on $\partial\Etil$ --- sums every fjord and tentacle
automatically, and yields the two-sided trapping with an \emph{absolute}
constant. Throughout, $\Etil$ denotes the main component
(Remark~\ref{rem:components}), $z$ and $\hat r=\sqrt{\hat\mu/\pi}$ are from
Corollary~\ref{cor:L1}, and we use the \emph{radial} and \emph{angular}
parts of the perimeter,
\[
   P_r:=\int_{\partial\Etil}|\nu\cdot e_r|\,d\HH^1,
   \qquad
   P_\theta:=\int_{\partial\Etil}|\nu\cdot e_\theta|\,d\HH^1,
   \qquad e_r=\tfrac{x-z}{|x-z|},
\]
where $\nu$ is the outward unit normal and $e_\theta\perp e_r$. Set
\[
   \rho_i:=\min_{x\in\partial\Etil}|x-z|,\qquad
   \rho_e:=\max_{x\in\partial\Etil}|x-z|
\]
(attained: $\partial\Etil$ is compact). Since $B_{\rho_i}(z)$ contains no
boundary point and surrounds $z\in\Etil$, $B_{\rho_i}(z)\subset\Etil$.

\begin{lemma}[Coarea identity for the angular perimeter]\label{lem:Ptheta}
With $N(\rho):=\HH^0\bigl(\partial\Etil\cap\{|x-z|=\rho\}\bigr)\in\{0,1,
2,\dots,\infty\}$ the number of boundary points at radius $\rho$,
\begin{equation}\label{eq:Ptheta-coarea}
   P_\theta=\int_0^\infty N(\rho)\,d\rho .
\end{equation}
\end{lemma}

\begin{proof}
The radial function $f(x):=|x-z|$ is $1$-Lipschitz; its tangential gradient
along the rectifiable curve $\partial\Etil$ has modulus $|\nabla^{\partial
\Etil}f|=|e_r-(e_r\cdot\nu)\nu|=\sqrt{1-(e_r\cdot\nu)^2}=|e_\theta\cdot\nu|$.
The coarea formula for Lipschitz functions on $1$-rectifiable sets
\textup{(}\cite[Theorem~2.93]{AFP2000}; \cite[\S3.4.2,
Theorem~3.13]{EvansGariepy2015}\textup{)} gives
$\int_{\partial\Etil}|\nabla^{\partial\Etil}f|\,d\HH^1
=\int_0^\infty\HH^0(\partial\Etil\cap f^{-1}(\rho))\,d\rho$, which is
\eqref{eq:Ptheta-coarea}.
\end{proof}

\begin{lemma}[Two crossings at every intermediate radius]\label{lem:Ncross}
If $\Etil$ is connected and every connected component of
$\Om\setminus\overline\Etil$ meets $\partial\Om$
\textup{(}Proposition~\ref{prop:noholes}\textup{)}, then $N(\rho)\ge2$ for
every $\rho\in(\rho_i,\rho_e)$.
\end{lemma}

\begin{proof}
Fix $\rho\in(\rho_i,\rho_e)$ and let $C_\rho:=\{|x-z|=\rho\}$, a circle. It
suffices to show $C_\rho$ meets both $\Etil$ and its complement, for then
$C_\rho$ (a connected curve) is split by $\partial\Etil$ into at least two
arcs, so $C_\rho\cap\partial\Etil$ has at least two points.

\emph{$C_\rho$ meets $\Etil$.} If $\rho\le\hat r$ this is clear, since
$B_{\rho_i}(z)\subset\Etil$ and $\Etil$ is connected with a point at radius
$\rho_e>\rho$, so by the intermediate value theorem applied to $|\cdot-z|$
on the connected set $\Etil$, $\Etil$ meets $C_\rho$. The same argument with
the point at radius $\rho_i\le\rho$ covers $\rho\ge\hat r$.

\emph{$C_\rho$ meets the complement.} Case $\rho>\rho_i$ \emph{and}
$\rho<\hat r$: the point $y\in\partial\Etil$ with $|y-z|=\rho_i$ is the tip
of a complement component $K$ (a connected component of
$\Om\setminus\overline\Etil$ adjacent to $y$); by
Proposition~\ref{prop:noholes}, $\overline K$ meets $\partial\Om$, and since
$\Etil\subset\subset\Om$ with $\overline\Etil\subset B_{\hat r+1}$ say while
$\partial\Om\not\subset B_{\hat r}$ (as $|\Om|=\pi>\hat\mu=|\Etil|$), $K$
contains points of radius $>\hat r>\rho$ as well as points of radius
arbitrarily close to $\rho_i<\rho$; being connected, $K$ meets $C_\rho$.
Case $\rho\ge\hat r$: since $|\Etil|=\pi\hat r^2<\pi\rho^2=|B_\rho(z)|$, the
disk $B_\rho(z)$ is not contained in $\Etil$, so $C_\rho$ (indeed
$B_\rho(z)$) meets the complement. In all cases $C_\rho$ meets both phases.
\end{proof}

\begin{theorem}[Two-sided trapping by the perimeter, no covering]
\label{thm:perim-annulus}
Let $\dT\le\delta_\star$ and let $\Etil$ be connected
\textup{(}Remark~\ref{rem:components}\textup{)}. Then
\begin{equation}\label{eq:perim-annulus}
   \rho_e-\rho_i\ \le\ \tfrac12\,P_\theta\
   \le\ \tfrac12\sqrt{2\,P(\Etil)\,(P(\Etil)-P_r)}\
   \le\ C\,\bigl(\mathrm{Exc}+a\bigr)^{1/2}\ \le\ C\,\dT^{1/4},
\end{equation}
where $\mathrm{Exc}:=P(\Etil)-2\pi\hat r=2\pi\hat r\,\diso(\Etil)\le C\dT$
and $a:=|\Etil\triangle B_{\hat r}(z)|\le C\sqrt\dT$. In particular
\[
   B_{\rho_i}(z)\subset\Etil\subset B_{\rho_e}(z),
   \qquad \rho_e-\rho_i\le C\,\dT^{1/4},
\]
with $C$ \emph{absolute}, and the multi-fjord bookkeeping is not needed.
\end{theorem}

\begin{proof}
\emph{Step 1: $P_\theta\ge2(\rho_e-\rho_i)$.} By
Lemma~\ref{lem:Ncross}, $N(\rho)\ge2$ on $(\rho_i,\rho_e)$, so by
Lemma~\ref{lem:Ptheta},
$P_\theta=\int_0^\infty N\,d\rho\ge\int_{\rho_i}^{\rho_e}2\,d\rho
=2(\rho_e-\rho_i)$.

\emph{Step 2: $P_\theta^2\le 2P(P-P_r)$.} By Cauchy--Schwarz,
$P_\theta^2=\bigl(\int|\nu\cdot e_\theta|\bigr)^2\le P\int(\nu\cdot
e_\theta)^2=P\int\bigl(1-(\nu\cdot e_r)^2\bigr)=P\bigl(P-\int(\nu\cdot
e_r)^2\bigr)$, and $\int(\nu\cdot e_r)^2\ge(\int|\nu\cdot e_r|)^2/P=P_r^2/P$,
whence $P_\theta^2\le P(P-P_r^2/P)=P^2-P_r^2\le 2P(P-P_r)$.

\emph{Step 3: $P-P_r\le\mathrm{Exc}+Ca$.} Let $R(\theta):=\sup\{\rho:
z+\rho e^{i\theta}\in\Etil\}\in[\rho_i,\rho_e]$ be the outer radial profile.
The radial slicing/area formula gives
$P_r=\int_{S^1}\bigl(\sum_{\rho:\,z+\rho e^{i\theta}\in\partial\Etil}1\bigr)
\rho\,\cdots$; we use only the elementary lower bound $P_r\ge\int_{S^1}
R(\theta)\,d\theta$ (the outermost crossing contributes at least its arc).
Where $R(\theta)<\hat r$, the sector $\{R(\theta)<\rho<\hat r\}$ at angle
$\theta$ lies in $B_{\hat r}(z)\setminus\Etil$, and since $R\ge\rho_i\ge
\hat r-\dT^{1/4}\ge\hat r/2$,
\[
   a\ \ge\ |B_{\hat r}(z)\setminus\Etil|
   \ \ge\ \int_{\{R<\hat r\}}\!\!\int_R^{\hat r}\!\rho\,d\rho\,d\theta
   \ \ge\ \frac{\hat r}{2}\!\int_{S^1}(\hat r-R)_+\,d\theta,
\]
so $\int_{S^1}(\hat r-R)\,d\theta\le\int_{S^1}(\hat r-R)_+\,d\theta\le
2a/\hat r$. Hence $\int R\,d\theta\ge2\pi\hat r-2a/\hat r$ and
\[
   P-P_r\ \le\ P-\!\int R\,d\theta\ \le\ \bigl(P-2\pi\hat r\bigr)+\frac{2a}{\hat r}
   \ =\ \mathrm{Exc}+\frac{2a}{\hat r}.
\]

\emph{Step 4: assemble.} Combining Steps 1--3 and $P=2\pi\hat r+\mathrm{Exc}
\le C$,
\[
   2(\rho_e-\rho_i)\le P_\theta\le\sqrt{2P(\mathrm{Exc}+2a/\hat r)}
   \le C(\mathrm{Exc}+a)^{1/2}.
\]
Finally $\mathrm{Exc}=2\pi\hat r\,\diso(\Etil)\le C\dT$ by
\eqref{eq:diso-inset}, and $a\le C\sqrt\dT$ by Corollary~\ref{cor:L1}, so
$(\mathrm{Exc}+a)^{1/2}\le C(\dT+\sqrt\dT)^{1/2}\le C\dT^{1/4}$. The
inclusions are the definitions of $\rho_i,\rho_e$ together with
$B_{\rho_i}(z)\subset\Etil$.
\end{proof}

\begin{remark}[Why the covering vanished]\label{rem:covering-gone}
The bookkeeping difficulty of Remark~\ref{rem:upper-status} was: many
fjords, each charged to the global $\diso$ without double counting. The
coarea identity \eqref{eq:Ptheta-coarea} does the bookkeeping in one stroke
--- $N(\rho)$ counts \emph{all} crossings at radius $\rho$ simultaneously,
so $\int N\,d\rho$ is automatically additive over fjords with no risk of
double counting, and the lone Cauchy--Schwarz of Step~2 converts it into the
global deficit $P-P_r\le\mathrm{Exc}+Ca$. No pointwise upper density and no
covering are used; the only ingredients are the deficit $D_I$ (through
$\mathrm{Exc}$ and, via FMP, through $a$), connectedness of $\Etil$, and the
no-hole property (Proposition~\ref{prop:noholes}). Notably $D_H$ is not even
needed for the trapping.
\end{remark}

\subsection{The regularity-free annulus}

\begin{theorem}[$C^2$-free trapping]\label{thm:dens-annulus}
Let $\dT\le\delta_\star$ and let $\Etil$ \textup{(}its main component, cf.\
Remark~\ref{rem:components}\textup{)} satisfy the two-sided density
\eqref{eq:two-sided} with the absolute constant $c_0$ of
Propositions~\ref{prop:lower-dens}--\ref{prop:upper-dens} and
\[
   r_*:=\max\{r_*^{\mathrm{int}},r_*^{\mathrm{ext}}\}
   \ \le\ C\,\dT^{1/4}.
\]
Then
\begin{equation}\label{eq:final-annulus}
   B_{\hat r-\eta}(z)\ \subset\ \Etil\ \subset\ B_{\hat r+\eta}(z),
   \qquad
   \eta\ \le\ C\,\dT^{1/4},
\end{equation}
with $C$ \emph{absolute} \textup{(}depending only on $c_0$ and the
universal constants of the deficit identities\textup{)}, and in particular
\emph{not on the $C^2$ geometry of $\Etil$}.
\end{theorem}

\begin{proof}
Apply Lemma~\ref{lem:chunk} with $A=a\le C\sqrt\dT$ (Corollary~\ref{cor:L1})
and $r_*\le C\dT^{1/4}$: then
$\eta=\max\{2r_*,2\sqrt{a/(\pi c_0)}\}\le C\max\{\dT^{1/4},(\sqrt\dT)^{1/2}\}
=C\dT^{1/4}$. The density hypothesis \eqref{eq:two-sided} is
Propositions~\ref{prop:lower-dens} (lower, $r\ge r_*^{\mathrm{int}}\le
C\dT$) and~\ref{prop:upper-dens} (upper, $r\ge r_*^{\mathrm{ext}}\le
C\dT^{1/4}$), valid for $r\in[r_*,r_0]$; the constants are absolute by those
propositions.
\end{proof}

\begin{remark}[What this achieves, and the loop back to $D_H,D_I$]
\label{rem:loop}
Theorem~\ref{thm:dens-annulus} replaces the $C^2$-dependent constant of
Theorem~\ref{thm:annulus} by an \emph{absolute} one: the only inputs are
the two-sided density and the $L^1$-closeness, and \emph{both reduce to the
deficits} --- the lower density to $D_H$ (the weighted Serrin defect $W\le
\tfrac{C_0}{4\pi}\dT$, Proposition~\ref{prop:lower-dens}), the upper density
and the $L^1$ scale to $D_I$ (the isoperimetric defect, FMP, and the fjord
perimeter cost, Propositions~\ref{prop:upper-dens} and
Corollary~\ref{cor:L1}). The annulus width $\eta\le C\dT^{1/4}$ is governed
by the exterior (fjord) scale $r_*^{\mathrm{ext}}\sim\diso^{1/2}$ and the
$L^1$ scale $\sqrt a\sim\dT^{1/4}$. With the annulus now in hand
\emph{without} regularity, one may feed it back into the level-set identity:
the level sets of $u$ inside the annulus inherit the trapping, and a second
application of $\int(D_H+D_I)=4\pi\dT$ on the trapped configuration is the
intended mechanism for improving the exponent --- the program's next step,
not carried out here. The multi-fjord covering bookkeeping that was the lone
residual is \emph{removed} by Theorem~\ref{thm:perim-annulus}: the
angular-perimeter coarea identity sums all fjords at once, so the two-sided
$C^2$-free annulus $\rho_e-\rho_i\le C\dT^{1/4}$ holds with absolute
constants, conditioned only on connectedness of the main component (which
carries area $\hat\mu-O(\dT)$, Remark~\ref{rem:components}).
\end{remark}
